%% en/sections/01_storica.tex
\lettrine[lines=4]{T}{o} understand the Sicilian Vespers, we must set
aside the nationalist reading that has dominated its collective memory
— that of «the Sicilian people rising up against the foreigner» — and
see Sicily in 1282 for what it was: a crucial node in a network of
geopolitical conflicts spanning the entire Mediterranean basin.

\nota[GM Note]{
  The «Aragonese conspiracy» orchestrated by Giovanni da Procida is a
  well-established but historically contested narrative. Historians such as
  Steven Runciman (\textit{The Sicilian Vespers}, 1958) and David Abulafia
  accept it in broad strokes, while others, including Jean Dunbabin,
  stress that the uprising of 30 March was largely spontaneous, fuelled by
  deep resentment and the accumulated weight of daily abuses.
  \textit{Vespri 1282} deliberately presents both readings as overlapping
  layers of reality: the conspiracy exists, but it does not control
  everything. It is this tension that makes the story playable.
}

\section{The Mediterranean in the 13th Century: A Contested Sea}

Sicily had never been simply «an island.» From the ninth to the eleventh
century it was an Islamic emirate of extraordinary cultural complexity.
Under the Norman conquest (1061–1091) it became the centre of a kingdom
that held together Arab, Greek, Latin and Norman cultures in a precarious
but creative equilibrium. Frederick II of Hohenstaufen (1194–1250) made
this kingdom a unique political and intellectual laboratory in Europe,
transforming Palermo into one of the most sophisticated capitals in the
medieval world.

Frederick II's death left a vacuum. His heirs — Conrad IV, then the young
Conradin — failed to hold power. Pope Innocent IV, who had excommunicated
Frederick and fought the Hohenstaufen for decades, had no intention of
allowing the Kingdom of Sicily to remain in the hands of that dynasty.
The papal solution was radical: offer the kingdom to a foreign royal house,
powerful enough to defeat the Hohenstaufen but dependent enough on Rome
to remain controllable.

The choice fell on Charles of Anjou, brother of Louis IX of France.

\secsep

\section{Charles of Anjou: A Royal Entrepreneur}

Charles of Anjou was not simply a conqueror. He was something more modern
and more dangerous: a political entrepreneur with imperial visions. He had
defeated Manfred of Hohenstaufen at the Battle of Benevento (1266) and
had the young Conradin beheaded in the Market Square of Naples (1268),
closing the Hohenstaufen dynasty with a calculated cruelty that shocked
all of Europe.

Established as king of Sicily and Naples, his project was the construction
of a Mediterranean empire with projections toward Constantinople and the
Levant. In 1267 he signed the Treaty of Viterbo with Baldwin II, the
Latin Emperor of Constantinople in exile. In 1277 he purchased the title
of King of Jerusalem. In early 1282, he was preparing the largest military
expedition of his career: a fleet of 300 ships aimed at recapturing
Constantinople.

\citastorica{%
  Charles had made the Mediterranean a personal matter.
  Sicily was his logistics base, his granary, his arsenal.
  The Sicilians were his instruments.%
}{Steven Runciman, \textit{The Sicilian Vespers}, 1958}

The problem was that the Sicilians — barons, clergy, craftsmen, farmers —
had no interest in being anyone's instruments.

\subsection{The Angevin Burden on Sicily}

Angevin rule translated into a system of taxation and military requisitions
unprecedented in the island's history. To finance his campaigns, Charles
imposed extraordinary taxes at a frequency that made any economic planning
impossible for the local population. French officials — the so-called
\textit{Provençals} or \textit{French}, a term Sicilians used
interchangeably for all continentals in Charles's service — were perceived
as hostile foreign bodies, often arrogant, sometimes violent.

The requisitions of ships, horses and provisions for the Byzantine
expedition — planned for the spring of 1282 — had reached unsustainable
levels. Everything weighed on Sicily.

\secsep

\section{The Conspiracy: Reality, Legend, Narrative Tool}

\subsection{Giovanni da Procida and the Aragonese Network}

Giovanni da Procida (c.~1210–1298) is the protagonist of the
«conspirationist» version of the Vespers. Physician, jurist, diplomat, he
had been one of the most loyal advisors of Frederick II and then of
Manfred. After the fall of the Hohenstaufen he had found refuge at the
court of Peter III of Aragon, who had married Constance, Manfred's
daughter, and believed he had legitimate rights to the Kingdom of Sicily.

According to historiographical tradition, da Procida had woven a network
connecting the Aragonese court of Barcelona, Pope Nicholas III (hostile to
Charles), the Byzantine emperor Michael VIII Palaeologus (who had every
interest in blocking Charles's expedition against Constantinople), and
Sicilian barons.

\nota[Historiographical Debate — Use in Play]{
  Modern historiography is divided. Runciman accepts the existence of a
  structured conspiracy but emphasises that the conspirators did not
  control the timing: the uprising of 30 March broke out earlier than
  planned, triggered by an incident no one had planned. Abulafia tends to
  reduce the weight of the conspiracy and amplify that of spontaneous
  popular resentment. \textbf{In the game, this ambiguity is a resource}:
  the Conspirators faction has a plan, but the plan is imperfect, and the
  other tables will act with or without them.
}

\subsection{The Role of the Church}

Pope Martin IV (elected 1281) was French, a creature of Charles, who had
excommunicated Michael VIII for political rather than theological reasons.
But the Roman curia contained currents hostile to Angevin hegemony.
Sicilian bishops, Italian cardinals, friars who had seen their communities
crushed by requisitions — the Sicilian clergy occupied a position of
extraordinary ambiguity. Churches were places of refuge and clandestine
communication. Friars moved freely between cities. Bishops had direct
contacts with Rome that bypassed Angevin administration. They were,
potentially, an already-existing intelligence network — but its members
owed obedience to Rome, not to Aragon.

\secsep

\section{Palermo in February 1282}

\subsection{The City and its Spaces}

Palermo in 1282 was a city of approximately 50,000 inhabitants, one of
the largest in Europe. Its layout bore the overlapping marks of seven
centuries of different dominations: the Arab city with its \textit{hara}
(neighbourhoods), the Norman city with its golden churches, the suburbs
that had grown around the gates.

\begin{multicols}{2}
\noindent\textbf{Main Quarters:}
\begin{itemize}[leftmargin=*, nosep]
  \item \textit{Cassaro} — the historic centre, seat of the royal palace
  \item \textit{Kalsa} — the Arab quarter, near the port
  \item \textit{Seralcadio} — craftsmen and merchants
  \item \textit{Albergheria} — the poorest, most densely populated quarter
\end{itemize}

\columnbreak

\noindent\textbf{Key Locations:}
\begin{itemize}[leftmargin=*, nosep]
  \item \textit{Santo Spirito} — the spark of the Vespers
  \item \textit{Palazzo dei Normanni} — seat of Angevin power
  \item \textit{Port of Palermo} — Charles's fleet in preparation
  \item \textit{Cathedral} — seat of the Archbishop
\end{itemize}
\end{multicols}

\subsection{The Tension in the Air}

The weeks before 30 March were marked by a growing tension that anyone
living in Palermo could feel, even without knowing anything of Giovanni
da Procida or Peter of Aragon. French soldiers had become more aggressive
in their searches. Rumours of new requisitions for the spring expedition
circulated in the markets. Someone said Charles was about to impose a new
census — which meant new taxes.

\secsep

\section{30 March 1282: The Beginning of the End}

\subsection{The Vespers at Santo Spirito}

On the evening of Easter Monday — 30 March 1282 — the Palermitans
gathered at the church of Santo Spirito for the festive procession.
What exactly happened — who started it, whether there was a deliberate
provocation or a spontaneous incident — is precisely the kind of historical
question that cannot be resolved with certainty. Sources agree on the
episode of French soldiers searching women in the crowd, interpreted as
violence and outrage; they agree on the sudden explosion of violence;
they agree that within a few hours the uprising had spread to the entire
city.

\citastorica{%
  All the people of the city, as though awaiting only that signal,
  rose against the French with knives and makeshift weapons, crying:
  \textit{Mora, mora li Francisi} — Death, death to the French.%
}{Giovanni Villani, \textit{Nuova Cronica}, 14th century}

By the morning of 31 March, almost all the French in Palermo —
approximately 2,000 people, soldiers, officials and civilians — had been
killed. The uprising spread rapidly to other Sicilian cities.
Charles of Anjou's imperial project was over.

\nota[Design Note — Why Stop at 30 March]{
  \textit{Vespri 1282} focuses on the \textbf{eve} of the uprising,
  not on the uprising itself. 30 March is the climax, but the story
  that matters to players is what precedes it: the choices, the alliances,
  the betrayals, the plans that work or collapse. The violence of the
  Vespers night is off-screen — it is what everything tends toward,
  not what is played.
}